\label{chap:pt1_conclusion}

This first part of the thesis focused on the development of an Intelligent Decision Support System (IDSS) for the fixture layout optimization problem in wood industry.

Several solution strategies were developed and evaluated. Initially, a simplified version of the problem was considered, which did not account for the presence of areas to be extruded within the workpiece and for the rotation of the fixtures. This problem was formulated as a CP model implemented in MiniZinc and evaluated using two alternative objective functions. The first, $\finer$, accounts for the principal moments of inertia of the fixture system, whereas the second, $\fdistsemi$, provides an alternative formulation that approximates $\finer$ by considering the distance and size of the fixtures.

The analysis conducted on a case study showed that $\finer$ is more suitable for identifying optimal fixture layouts, but its computational cost increases substantially with problem complexity. In contrast, $\fdistsemi$ provides good-quality solutions in shorter computation times and is less affected by increasing problem complexity. These results therefore motivated the use of an alternative objective function that approximates the results produced by $\finer$ for exact optimization techniques such as CP or MIP.

The complete problem was subsequently considered, incorporating both the areas to be extruded and the rotation of the fixtures. The CP model was extended and compared with two alternative approaches: a MIP model and a Q-learning based algorithm. Testing the three approaches on real-world case studies, the CP model proved to be the most effective solution, with its effectiveness demonstrated by its ability to obtain optimal solutions for the simpler instances when solved with Chuffed, while the use of $\lns$ enabled it to handle the more complex cases. The MIP model solved with Gurobi produced solutions comparable to those of CP for simpler instances, but its performance deteriorated as the problem complexity increased. Q-learning was able to generate good fixture layouts, but did not outperform the solutions obtained with CP or MIP.

None of these approaches, however, explicitly considered fixture rotation during the optimization process, as including arbitrary rotations directly in the decision variables would substantially increase the search space. Rotation was therefore introduced in a subsequent refinement phase based on the PSO meta-heuristic. Starting from the solutions generated by the three approaches, PSO improved the initial layouts in all test cases and also outperformed the layouts provided by an expert operator, with the exception of one case in which the initial solution was of low quality. These results demonstrate the effectiveness of using local search as a refinement mechanism and highlight the advantage of integrating the approaches, as the initial feasible layouts provide PSO with a suitable starting point without requiring it to search for feasibility from scratch.

The application of PSO to the fixture layout optimization problem relied on the linear constraints of the MIP model, implemented in MiniZinc. However, incorporating these constraints into PSO required them to be explicitly reformulated and integrated in its code. To overcome this limitation and enable PSO to solve MiniZinc models involving continuous variables and nonlinear constraints, a FlatZinc-based PSO algorithm was developed. The algorithm exploits the information contained in the FlatZinc model, obtained form the compilation of the MiniZinc model, and represents its information through an invariant graph, allowing constraint violations and the values of defined variables to be evaluated directly from this representation.

The comparison with a PSO implementation in which the variables and constraints are directly encoded in the meta-heuristic showed that the FlatZinc-based approach can achieve better solution quality while using the same search mechanism. This improvement can therefore be attributed to the additional information provided by the FlatZinc model, which offers a more explicit description of the problem and can be exploited to guide the search.

Overall, the results demonstrate that complex constrained problems arising in industrial applications often benefit from the integration of exact optimization and local search techniques. The results obtained for the fixture layout optimization problem demonstrate the effectiveness of this integration, while the development of the FlatZinc-based PSO extends this approach beyond the specific application considered and opens the possibility of applying PSO to a broader class of nonlinear models formulated in MiniZinc.


